\documentclass[10pt,a4paper]{amsart}
\usepackage[margin=19mm]{geometry}
\usepackage{amsmath,amssymb,mathtools}
\usepackage[scr=rsfs,cal=euler]{mathalfa}
\usepackage{xfrac}
\usepackage[unicode,hidelinks,bookmarks=false]{hyperref}
\newcommand{\ie}{i.e.\ }
\newcommand{\cf}{cf.\ }
\newcommand{\resp}{respectively\ }
\newcommand{\Subsection}{\subsection}
\providecommand{\nobreakdash}{\nobreak\hbox{-}}
\setlength{\emergencystretch}{2em}
\multlinegap0pt
\usepackage{bm}
\usepackage{mathtools}
\usepackage{xcolor}
\usepackage{float}
\usepackage[shortlabels]{enumitem}
\newtheorem{proposition}{Proposition}
\newcommand{\DDSM}{\text{DDSM}}
\newcommand{\R}{\mathbb{R}}
\def\rvx{{\mathbf{x}}}
\def\vs{{\bm{s}}}
\def\vtheta{{\bm{\theta}}}
\title{When Are Two Scores Better Than One? \\ Investigating Ensembles of Diffusion Models}
\author{Rapha\"el Razafindralambo, R\'emy Sun, Fr\'ed\'eric Precioso, Damien Garreau, Pierre-Alexandre Mattei}
\date{Source: arXiv:2601.11444v2 (CC BY 4.0)}
\begin{document}
\maketitle

\noindent\emph{Source and licence.} When Are Two Scores Better Than One? \\ Investigating Ensembles of Diffusion Models. Version arXiv:2601.11444v2 (CC BY 4.0), distributed by arXiv under the Creative Commons Attribution 4.0 International licence, http://creativecommons.org/licenses/by/4.0/, as recorded in the arXiv licence metadata for this exact version and shown on the abstract page https://arxiv.org/abs/2601.11444v2. Full licence notice: \texttt{licenses/arxiv-2601.11444-CC-BY-4.0.txt}.\par\smallskip

\noindent\textit{Editorial scope.} Counted unit: the article's proposition loss (source0.tex lines 1318--1327). The article's complete original proof of it is delivered with it (source0.tex lines 1328--1385, one \texttt{proof} environment of 58 lines). No mathematics is written, repaired or re-derived: the statement and the proof are the article's own lines, in the article's own notation, and no retained line is substituted or re-flowed.  No further source-local context is needed: every cross-reference of every retained line resolves inside the two delivered spans.  Nothing was omitted from any retained span, so the package carries no omission record.  No retained line contains a citation, so the wrapper carries no bibliography and no bibliography entry.  No retained line contains a figure environment, an image-inclusion command or drawing code, so no image is delivered and the package declares no asset dependency.  Editorial formatting only, in addition: an AMS \texttt{amsart} preamble that restates the article's own declarations and macros used by the retained lines, namely the environments \texttt{proposition} on separate counters, together with the standard packages those lines need.  The retained environments print in this document as Proposition 1 in order of appearance, whereas the article prints the same items with the numbering of its own full document.  The complete native source of the article as distributed in the arXiv e-print for this exact version is bundled unchanged as \texttt{sources/192866/source0.tex}.
\begin{proposition}\label{proposition: ensemble monotonicity score loss}
Let $\vs^{(1)}_\vtheta,\dots,\vs^{(K)}_\vtheta$ be $K$ score estimators mapping $\R^d \times [0,T]$ to $\R^d$. Then, the DDSM loss satisfies
\begin{equation}
\mathcal{L}_\DDSM(\overline{\vs}_\vtheta^{(K)},\lambda(\cdot)) \leq \frac{1}{K} \sum_{j=1}^K \mathcal{L}_\DDSM\left(\frac{1}{K-1} \sum_{k \neq j} \vs_\vtheta^{(k)},\:\:\lambda(\cdot)\right)
\end{equation}
Moreover, if for any $(\rvx_t,t) \in \R^d \times [0,T]$, the outputs $\vs^{(1)}_\vtheta(\rvx_t,t), \dots, \vs_\vtheta^{(K)}(\rvx_t,t) $ are identically distributed, then, 
\begin{equation}
\mathbb{E} \left[\mathcal{L}_{\DDSM}(\overline{\vs}_\vtheta^{(K+1)},\lambda(\cdot)) \right] \leq \mathbb{E}\left[\mathcal{L}_{\DDSM}(\overline{\vs}_\vtheta^{(K)},\lambda(\cdot)) \right].
\end{equation}
\end{proposition}

\begin{proof}
Let $K>1$. As a first step, we establish the following algebraic identity:
\begin{equation}
\frac{1}{K}\sum_{k=1}^K {\vs}_\vtheta^{(k)}(\rvx_t,t)= \frac{1}{K} \sum_{j=1}^K \frac{1}{K-1} \sum_{k \neq j} {\vs}_\vtheta^{(k)}(\rvx_t,t),
\end{equation}
which is a direct consequence of
\begin{equation}
\begin{bmatrix}
1 \\[2pt]
1 \\[2pt]
\vdots \\[2pt]
1 \\[2pt]
1
\end{bmatrix}
= \frac{1}{K-1}
\left(
\begin{bmatrix}
0 \\[2pt]
1 \\[2pt]
\vdots \\[2pt]
1 \\[2pt]
1
\end{bmatrix}
+
\begin{bmatrix}
1 \\[2pt]
0 \\[2pt]
\vdots \\[2pt]
1 \\[2pt]
1
\end{bmatrix}
+ \cdots +
\begin{bmatrix}
1 \\[2pt]
1 \\[2pt]
\vdots \\[2pt]
1 \\[2pt]
0
\end{bmatrix}
\right). 
\end{equation}
We now use Jensen's inequality and the convexity of the norm to obtain the first inequality.
\begin{align} 
\mathcal{L}_\DDSM(\overline{\vs}_\vtheta^{(K)},\lambda(\cdot)) & =  \frac{1}{2} \int_0^T  \mathbb{E}_{p_0(\rvx_0)p_{t|0}(\rvx_t|\rvx_0)}\left[\lambda(t) \| \left(\frac{1}{K}\sum_{k=1}^K\vs_\vtheta^{(k)}(\rvx_t,t)\right) - \nabla_{\rvx_t} \log q_{t|0} (\rvx_t | \rvx_0) \|^2_2\right] \\
& =  \frac{1}{2} \int_0^T  \mathbb{E}_{p_0(\rvx_0)p_{t|0}(\rvx_t|\rvx_0)}\left[\lambda(t) \| \frac{1}{K} \sum_{j=1}^K \frac{1}{K-1} \sum_{k \neq j} {\vs}_\vtheta^{(k)}(\rvx_t,t) - \nabla_{\rvx_t} \log q_{t|0} (\rvx_t | \rvx_0) \|^2_2\right] \\
& \leq \frac{1}{2K}\sum_{j=1}^K \int_0^T  \mathbb{E}_{p_0(\rvx_0)p_{t|0}(\rvx_t|\rvx_0)}\left[ \lambda(t) \| \frac{1}{K-1} \sum_{k \neq j} {\vs}_\vtheta^{(k)}(\rvx_t,t) - \nabla_{\rvx_t} \log q_{t|0} (\rvx_t | \rvx_0) \|^2_2\right] \\
& \eqqcolon \frac{1}{K} \sum_{j=1}^K \mathcal{L}_\DDSM\left(\frac{1}{K-1} \sum_{k \neq j} \vs_\vtheta^{(k)},\:\:\lambda(\cdot)\right).
\end{align}

Now assume that for any pair $(\rvx_t,t)$, the random variables $\vs^{(1)}(\rvx_t,t),\dots,\vs^{(K)}(\rvx_t,t)$ are identically distributed given $(\rvx,t)$. Then,

\begin{align}
\mathbb{E}[\mathcal{L}_\DDSM(\overline{\vs}_\vtheta^{(K)},\lambda(\cdot))] & = \frac{1}{K} \sum_{j=1}^K \mathbb{E}  \left[ \mathcal{L}_\DDSM\left(\frac{1}{K-1} \sum_{k \neq j} \vs_\vtheta^{(k)},\:\:\lambda(\cdot)\right) \right] \\
& = \frac{1}{K} \sum_{j=1}^K \mathbb{E}  \left[ \mathcal{L}_\DDSM\left(\frac{1}{K-1} \sum_{k = 1}^{K-1} \vs_\vtheta^{(k)},\:\:\lambda(\cdot)\right) \right] \\
& = \mathbb{E}  \left[ \mathcal{L}_\DDSM\left(\overline{\vs}_\vtheta^{(K-1)},\lambda(\cdot)\right) \right].
\end{align}

\end{proof}

\end{document}
